% Entwurf Abschnitt 6.4 (AP6), Stand 20.11.2026. Budget ca. 2,5 Seiten (0,8 / 1 / 0,7).
% Code-Stand aller Befunde: G1@25d5c1a. Pfade in % Basis: relativ zu g1pilot/ bzw. Repo-Root.
% Darstellungen (baut ein anderer Auftrag, hier nur referenziert): tab:bedienwege, fig:posen_ablaeufe.
\section{Bedienkonzept}
\label{sec:bedienkonzept}
Das Bedienkonzept setzt die Anforderung~A6 um, die Umgebung ohne Vorerfahrung zu bedienen, und verbindet sie mit
dem sicheren Realbetrieb nach A5 (Tabelle~\ref{tab:anforderungen}).
Auf fünf Leitgedanken (Abschnitt~\ref{sec:leitgedanken}) folgen die Bedienwege und ihr Zusammenwirken
(Abschnitt~\ref{sec:bedienwege_konzept}) sowie das Einrichten von Aufgaben ohne Programmierung
(Abschnitt~\ref{sec:posen_ablaeufe}); einzelne Bedienschritte stehen in Anhang~\ref{anh:dokumentation}.

\subsection{Leitgedanken}
\label{sec:leitgedanken}
Die Zielgruppe nach A6 sind Personen ohne Erfahrung in Robotik und Programmierung.
Für sie passen die Teleoperationssysteme aus Abschnitt~\ref{sec:teleoperation_teil} nur bedingt, weil
Motion-Capture-Systeme, VR-Brillen und Exoskelette teuer und an bestimmte Einsatzorte gebunden sind
\cite{Fu2024}\ev{Fu2024.01} und das System des Herstellers vorab das Lesen der offiziellen Dokumentation verlangt
\cite{UnitreeXRTeleop}\ev{UnitreeXRTeleop.02}.
Die Umgebung setzt deshalb auf eine Bedienoberfläche am PC und Standard-Eingabegeräte, deren Gestaltung fünf
Leitgedanken folgt.

% Basis: G1@25d5c1a docs/10-A-Demo-GUI.md „Ziel“, „1 · Modus (oben)“, „2 · Steuerung (Mitte)“
Der erste Leitgedanke verlangt, dass die bedienende Person zu jedem Zeitpunkt nur eine Entscheidung trifft.
Von den zwei Bedienmodi, GEHEN für die Fortbewegung und GREIFEN für Arme und Hände, ist genau einer aktiv, und
angezeigt wird nur, was im aktiven Modus sinnvoll ist.
Eine Oberfläche mit allen Funktionen zugleich eignet sich für die Entwicklung, lässt aber offen, in welchem Zustand
der Roboter ist und welche Eingabe gerade wirkt.
\qf{Beleg zu Modusfehlern bzw. reduzierter Auswahl in Bedienoberflächen (Lücke Q4)}

% Basis: G1@25d5c1a docs/10-A-Demo-GUI.md „1 · Modus (oben)“ (Zustände der Kacheln, Startleiste, Bodenkontakt);
% docs/03-A-Startmenue.md „Optionen“ (Zeile Sicherheit); docs/04-T-Umgebungsvariablen.md (G1_REAL_CONFIRM)
Der zweite Leitgedanke macht den Zustand sichtbar, bevor eine Eingabe wirkt.
Jeder Bedienmodus erscheint unterscheidbar als aktiv, wechselnd oder inaktiv, und im Realbetrieb zeigt die
Bedienoberfläche den gemeldeten Zustand des Onboard-Reglers und das Ergebnis der Bodenkontaktprüfung, bevor das
Gehen freigegeben wird.
Nach dem dritten Leitgedanken „Simulation zuerst“ sind Bedienoberfläche und Bedienwege in beiden Welten dieselben
(Abschnitt~\ref{sec:sim_real_schalter}), sodass jede Bedienung zuerst gefahrlos geübt wird; der Start des
Realbetriebs verlangt zusätzlich die Bestätigung der Sicherheitshinweise.

% Basis: G1@25d5c1a docs/51-A-Notaus.md „Stufen“, „Was zusätzlich von selbst anhält“, „Altes Verhalten zurückholen“;
% docs/12-A-Handy.md (Hinweis Handy-NOT-HALT über den PC, Rollen); docs/50-A-Echter-Roboter.md E („Kein Auto-Start“);
% docs/10-A-Demo-GUI.md „Seite ARME“ (Abbruch eines Ablaufs); docs/40-T-Simulation.md „Auto-Start (GREIFEN)“
Der vierte Leitgedanke hält den Halt immer erreichbar: Der Not-Halt liegt auf Bedienoberfläche, Gamepad und
Smartphone und wirkt unabhängig vom Bedienmodus, von einer Sperre des Gamepads und in der Grundeinstellung auch aus
der Rolle, die das Smartphone nur zur Ansicht berechtigt.
Weil der Not-Halt des Smartphones über den PC läuft, bleibt die Fernbedienung des Herstellers die vom PC
unabhängige Halt-Stufe (Abschnitt~\ref{sec:zustaende}).
Nach dem fünften Leitgedanken bewegt sich im Realbetrieb nichts von selbst: Ohne Eingabe sendet die Umgebung keinen
Bewegungsbefehl, nach einem quittierten Not-Halt bleiben die Arme schlaff, bis sie bewusst übernommen werden, und
ein abgebrochener Ablauf setzt sich nie selbsttätig fort.
Ob die Leitgedanken die Bedienung ohne Vorerfahrung tragen, untersucht Abschnitt~\ref{sec:eval_bedienaufwand}.

\subsection{Bedienwege und geteilter Bedienzustand}
\label{sec:bedienwege_konzept}
% Basis: G1@25d5c1a docs/10-T-Demo-GUI.md „Grundsatz“; docs/12-A-Handy.md (Kopf: „Der PC rechnet weiter alles“)
Mehrere Bedienwege bedienen unterschiedliche Personen und Situationen, ohne die Logik zu vervielfachen.
Die Schnittstellen sind das Mittel der Interaktion zwischen Nutzer und Roboter \cite{Darvish2023}\ev{Darvish2023.04};
statt einer einzigen bietet die Umgebung mehrere an, die nach Abschnitt~\ref{sec:komponenten} keine Bewegung
berechnen, sondern Ziele an die Funktionsknoten senden.
Dadurch gelten die Schutzmechanismen aus Abschnitt~\ref{sec:schutzmechanismen} für jeden Bedienweg gleich.
Tabelle~\ref{tab:bedienwege} ordnet die Bedienwege den Funktionen zu.

% Basis: G1@25d5c1a docs/10-A-Demo-GUI.md; docs/11-A-Controller.md (Kopf, „Echter Roboter“); docs/12-A-Handy.md
% (Kopf, „Was das Handy nicht kann“, Hinweis „Zusammenstellen geht nur am PC“); docs/21-A-Arm-API.md „Überblick“;
% docs/02-T-Architektur.md „Projektstruktur — Node-Übersicht“ (interactive_marker)
Die Bedienoberfläche deckt alle Funktionen ab und ist der einzige Ort, an dem Abläufe zusammengestellt werden.
Das Gamepad dient der unmittelbaren Führung der Hände und dem Gehen mit Totmann-Prinzip, während das Smartphone den
Zustand der Bedienoberfläche spiegelt und deren Funktionen mit denselben Prüfungen aufruft, Gehen und unmittelbare
Führung der Hände aber bewusst nicht anbietet.
Zielmarker in RViz und die Programmierschnittstelle der Arme richten sich an Entwickler und externe Programme
(Abschnitt~\ref{sec:andockpunkte}); alle Bedienwege stehen in Simulation und Realbetrieb zur Verfügung.

% Basis: G1@25d5c1a g1pilot/teleoperation/shared_state.py (Modulkommentar); docs/50-A-Echter-Roboter.md E
% („Ein Ziel ersetzt das andere“); docs/23-A-Posen-und-Ablaeufe.md „Pose anfahren“ (manueller Vorrang)
Gleichzeitig aktive Bedienwege bergen widersprüchliche Anzeigen und konkurrierende Befehle.
Bedienoberfläche und Gamepad teilen sich deshalb einen Bedienzustand aus Bedienmodus, Auswahl der Hand und
Tempostufen, in dem je Feld die jüngste Änderung gilt; Fahrbefehle enthält er nicht, und jede Tempostufe begrenzt
jeder Bedienweg auf sein eigenes Maximum.
Für Bewegungsziele ersetzt ein neues Ziel das vorige, gleich von welchem Bedienweg es stammt, und eine manuelle
Eingabe bricht eine laufende geplante Bewegung ab.

% Basis: G1@25d5c1a docs/12-T-Handy.md „Aktionen und Rechte“; g1pilot/teleoperation/phone/gateway.py (Modulkommentar,
% ACTIONS); docs/12-A-Handy.md „Rollen“
Für das drahtlos angebundene Smartphone kommen Rollen und eine Steuerhoheit hinzu.
Die Rolle „Ansicht“ zeigt nur den Zustand, die Rolle „Bediener“ sendet Befehle, sobald das Gerät die Steuerhoheit
übernommen hat, die die Bedienoberfläche jederzeit zurückholen kann.
Es besteht nur eine Sitzung, und ein bewegender Befehl gilt nur mit Bezug auf einen aktuellen Zustand, sodass ein
verspätet eintreffender Befehl wirkungslos bleibt.

% Basis: G1@25d5c1a docs/51-A-Notaus.md „Quittieren“; docs/11-A-Controller.md „NOT-HALT und Quittieren“
Zwischen Auslösen und Aufheben eines Not-Halts besteht eine bewusste Asymmetrie: Ausgelöst wird er an Bedienoberfläche,
Gamepad und Smartphone, quittiert nur an Bedienoberfläche oder Smartphone, die den Zustand anzeigen, am Gamepad
dagegen nicht.
\pruefen{Begründung der Asymmetrie bestätigen; die Doku nennt nur die Doppelbelegung der Start-Taste am Gamepad}
Die Anbindung der Bedienwege beschreibt Abschnitt~\ref{sec:bedienwege}, den Bedienaufwand bewertet
Abschnitt~\ref{sec:eval_bedienaufwand}.

\subsection{Posen und Abläufe als aufgabenunabhängige Bausteine}
\label{sec:posen_ablaeufe}
% Basis (Alternativen und Entscheidung): Plan Kap. 6, Abschnitt 5, 6.4.3; Umsetzung G1@25d5c1a docs/23-T-Posen-und-Ablaeufe.md
Nach A1 soll sich eine neue Aufgabe ohne Programmierung einrichten lassen.
Ein Skript scheidet aus, weil Roboterprogrammierung neben Programmierkenntnissen Wissen aus mehreren Disziplinen
erfordert \cite{Ajaykumar2021}\ev{Ajaykumar2021.02}; weitere Alternativen sind grafische Umgebungen wie Choregraphe,
die ohne Kompilieren von C++ auskommen \cite{Pot2009}\ev{Pot2009.06}, und Lernen aus Demonstrationen auf Basis von
Teleoperationsdaten \cite{Cheng2024}\ev{Cheng2024.01}.
Die Umgebung beschreibt eine Aufgabe stattdessen als Daten aus wenigen festen Bausteinen, ähnlich der Bibliothek
vorprogrammierter Verhalten in Choregraphe \cite{Pot2009}\ev{Pot2009.05}.
\pruefen{Gründe gegen grafische Programmierung und Lernen aus Demonstrationen von Nils}

% Basis: G1@25d5c1a docs/23-A-Posen-und-Ablaeufe.md „Überblick“, „Pose speichern“; docs/23-T-Posen-und-Ablaeufe.md
% „Dateiformat“ (moves: je Achse begrenzt, Pause je Schritt); g1pilot/manipulation/pose_store.py (MAX_MOVE_M, MAX_WAIT_S)
Abbildung~\ref{fig:posen_ablaeufe} zeigt die drei Bausteine des Datenmodells.
Eine Pose ist eine benannte Stellung von Armen und Händen, eine relative Bewegung verschiebt die Hände von ihrer
aktuellen Stelle um einen je Achse begrenzten Betrag, und ein Ablauf ist eine Folge von Verweisen auf beide mit
optionaler Pause, sodass dieselbe Pose in mehreren Abläufen stehen kann.
Eingelernt wird eine Pose, indem Arme und Hände über einen Bedienweg in Stellung gefahren und gespeichert werden,
anders als bei der Handführung, in der der Bediener den Roboter dafür direkt bewegt \cite{Villani2018}\ev{Villani2018.06}.

% Basis: G1@25d5c1a docs/23-A-Posen-und-Ablaeufe.md „Ablauf abfahren“; docs/23-T-Posen-und-Ablaeufe.md „Dateiformat“,
% „Speicherort“; g1pilot/teleoperation/demo_gui.py default_sequence_mode() (Sim fortlaufend, real schrittweise)
Abläufe laufen fortlaufend oder schrittweise mit Bestätigung je Schritt, im Realbetrieb ist schrittweise
voreingestellt.
Jeder Schritt mit Armbewegung wird wie eine einzelne Zielvorgabe geplant und auf Kollisionen geprüft; der Ablauf
selbst enthält keine Bewegungslogik und umgeht keinen Schutzmechanismus.
Weil die Abläufe außerhalb der Container liegen und in beiden Welten dieselben sind, wird eine Aufgabe in der
Simulation eingerichtet und am Roboter ausgeführt.
% Kann
Offline-Programmierung an einer virtuellen Simulation \cite{Ajaykumar2021}\ev{Ajaykumar2021.07} kann Nutzeraufwand,
Kosten und Wartezeiten auf einen verfügbaren Roboter verringern \cite{Ajaykumar2021}\ev{Ajaykumar2021.08}.

% Basis: G1@25d5c1a docs/23-T-Posen-und-Ablaeufe.md (keine Bedingungen im Dateiformat); docs/23-A-Posen-und-Ablaeufe.md
% „Fehlerbehebung“ (Roboter steht anders als beim Speichern), Beispiel-Abläufe „gelten genau an der jeweiligen Station“
Die Grenze des Konzepts sind feste Folgen ohne Verzweigungen, Bedingungen an Sensorwerte oder Objekterkennung.
Zudem gelten die Posen nur, solange der Roboter so zum Arbeitsplatz steht wie beim Einlernen.
Wie weit sich Aufgaben dennoch einrichten lassen, zeigt Abschnitt~\ref{sec:eval_aufgaben}; die Grenze greift
Abschnitt~\ref{sec:limitierungen} auf.
